\begin{table}[tb]
\centering\small\setlength{\tabcolsep}{3pt}
\caption{全８条件の対比較．既知は両実行の全使用量が既知の対数／３．要求・出力は一致対数／判定対数．時計は最初の要求差が現在の時計欄だけだった観測対数．}
\begin{tabular}{llrrrr}
課題 & 条件 & 既知 & 要求 & 公開出力 & 時計 \\ \hline
写真 & M & 3/3 & 0/3 & 0/3 & 3 \\
写真 & H & 3/3 & 3/3 & 3/3 & 0 \\
精算 & M & 3/3 & 0/3 & 0/3 & 3 \\
精算 & H & 3/3 & 3/3 & 3/3 & 0 \\
返金 & M & 3/3 & 0/3 & 0/3 & 3 \\
返金 & H & 3/3 & 3/3 & 3/3 & 0 \\
書出し & M & 3/3 & 0/3 & 0/3 & 3 \\
書出し & H & 0/3 & 未確定 & 未確定 & 0 \\
\hline
\end{tabular}
\end{table}
要求列または公開出力列の完全一致は，全使用量が既知で比較できた対について示す．不明生成を含む対の全体一致は未確定とし，判定できた接頭部だけを全体へ広げない．時計欄以外をそろえた照合は観測された最初の要求差だけに関する結果であり，出力変化や全仕事費用の原因を単独で決める証拠ではない．
